% =============================================================================
% reader_notes.tex -- reader-facing notes, no project-path references
% =============================================================================
\noindent The body has deliberately moved quickly through a lot of material:
axioms, algebraic constructions, a numerical match, and then a boundary where a
tempting derivation stops short. These notes slow down at the places a careful
reader is most likely to ask, ``Wait, what exactly is being claimed here?'' They
do not add new results; they unpack the rules of reading the results already
given.

\section*{How to read the tags}
\addcontentsline{toc}{section}{How to read the tags}
\markboth{Reader's Notes}{How to read the tags}

\noindent The epistemic tags are not decorative labels. They are the grammar of
the argument. A \etag{thm} tag says: if you accept the axioms and the cited
classical mathematics, the statement follows. A \etag{par} tag says: a standard
physics formula has been supplied with FTD inputs, so the numerical agreement is
a cross-check rather than a derivation. A \etag{smc} tag says: the evidence is
substantial enough to matter, but there is still no chain from the axioms to the
physical identification. An \etag{opn} tag says: the question is live, and the
work needed to settle it has been named.

\medskip
\noindent This is why the same page can say that a polynomial is theorem-grade
while its physical interpretation is not. The algebraic statement
``this quadratic has these roots'' is a different species of claim from the
physics statement ``the larger root is the measured inverse electromagnetic
coupling.'' The first is a calculation inside the construction. The second is a
bridge to nature, and it has to earn a stronger readout rule before it can be
called a derivation.

\section*{Why the alpha match is not a derivation}
\addcontentsline{toc}{section}{Why the alpha match is not a derivation}
\markboth{Reader's Notes}{Why the alpha match is not a derivation}

\noindent The large root $x_+$ lands very close to $\alpha^{-1}$, and the match
is good enough that dismissing it casually would be intellectually lazy. But a
match, even a striking one, is not the same thing as a mechanism. To derive the
fine-structure constant, the framework would need to show why the physical
electromagnetic readout must use the larger root of the master quadratic.

\medskip
\noindent The missing step is not a better numerical approximation. The missing
step is an operational rule. In the language of Part~II, the framework must
produce a rank-2 readout operator whose trace and determinant are exactly
\[
  (\Tr,\Det)=(16G^{*2},16G^{*3})
\]
for a structural reason internal to the substrate. The trace and the odd
$G^*$ factor both have native sources. What is not forced is their assembly into
the same operator, with the odd factor placed specifically in the determinant
slot. That placement is the difference between a discovered readout and a chosen
one.

\section*{What the square-root wall means}
\addcontentsline{toc}{section}{What the square-root wall means}
\markboth{Reader's Notes}{What the square-root wall means}

\noindent A quadratic has two roots, and the symmetric data of the quadratic do
not by themselves choose between them. The trace and determinant know the pair of
roots together; they do not tell us which root nature should read as
$1/\alpha$. The quantity
\[
  \sqrt{G^*(4G^*-1)}
\]
is the first ingredient that distinguishes the two branches. That is why the
boundary is described as a square-root wall: the framework can build the
symmetric side of the quadratic, but it does not yet natively produce the
antisymmetric branch-selection object.

\medskip
\noindent Put less technically: FTD can get to the menu of possible readings,
but it does not force the dish. The present theory allows the alpha-like reading
and also allows other determinant choices over the same trace. Because both
sorts of readout are coherent, the alpha reading is not forced by the current
postulates alone.

\section*{What would close the gap}
\addcontentsline{toc}{section}{What would close the gap}
\markboth{Reader's Notes}{What would close the gap}

\noindent There are two honest ways the gap could close. The first would be a
new binding law: a substrate-native rule that fixes the readout assembly rather
than choosing it after seeing the target. If such a law were added and justified,
the alpha match would become a derivation relative to that added law.

\medskip
\noindent The second would be a direct substrate measurement of the binding
phenomenon. That route has not produced a positive result in the tested regime.
A finite null is not a universal impossibility theorem, but it does mean the
empirical path has not supplied the missing readout either.

\medskip
\noindent So the present conclusion is deliberately modest. The monograph does
not say that no future extension could ever derive $\alpha$. It says that the
current five-postulate ontology plus the algebraic spine does not already contain
the binding law that would make the derivation go through.

\section*{How to read the physics layer}
\addcontentsline{toc}{section}{How to read the physics layer}
\markboth{Reader's Notes}{How to read the physics layer}

\noindent The broader physics catalogue should be read with the same discipline.
Dimensionless matches are the most informative, because they do not depend on a
choice of metres, kilograms, seconds, or calibration anchors. Dimensional
quantities are weaker: once a length scale or mass scale is declared, many
numbers inherit that declaration.

\medskip
\noindent Likewise, a particle mass or decay rate computed by inserting FTD
numbers into a familiar formula from quantum field theory is not automatically an
FTD derivation. It may still be useful. It may show that the framework's
integers and constants sit in the right numerical neighborhood. But the
functional form is borrowed, so the honest tag is \etag{par}. The point of the
tagging system is not to diminish those matches; it is to keep them in the
right drawer.
